\chapter{究竟要证明什么}
\label{chap:overview}
\section{把命题写准确}
当复数 $s$ 的实部大于 $1$ 时，黎曼 $\zeta$ 函数定义为
\[
\zeta(s)=\sum_{n=1}^{\infty}n^{-s}
=\prod_{p\ \text{为素数}}(1-p^{-s})^{-1}.
\]
这些级数与乘积在这里绝对收敛。随后把 $\zeta$ 亚纯延拓到复平面；$s=1$ 是单极点，不能作为普通有限函数值参与零点讨论。

若 $\zeta(\rho)=0$，则称 $\rho$ 为零点。负偶数
$-2,-4,-6,\ldots$ 是平凡零点。本稿的目标为
\begin{theorem}[黎曼猜想的精确陈述]
\label{thm:RH}
若 $\rho\in\mathbb C$，$\rho\ne1$，$\zeta(\rho)=0$，且
$\rho\notin\{-2,-4,-6,\ldots\}$，则
\[
\Rea\rho=\frac12.
\]
\end{theorem}
这个命题说的是所有符合条件的零点。计算出很多零点在临界线上，不能替代全称证明。

\section{证明只需要两条主线相遇}
第一条主线来自算术。它构造一个记录素数与素数幂的实函数 $G(x)$，并在任意远的某些完整区间上，证明
\[
G(x)\le x^{1/2+K\varepsilon}.
\]
这里 $K$ 先固定，$\varepsilon>0$ 后选择，可以任意小。

第二条主线来自复分析。它证明，只要任意远处都有完整区间满足
$G(x)\le x^\theta$，就不可能存在 $\Rea\rho>\theta$ 的 $\zeta$ 零点。

把两条主线接起来，就能把排除区域的边界推到 $1/2$。最后利用函数方程，临界线左侧的非平凡零点会对应到已经排除的右侧，因而也不可能存在。

\[
\boxed{\text{算术窗口}}
\ \Longrightarrow\ 
\boxed{\Rea\rho\le\tfrac12}
\ \xRightarrow{\text{函数方程}}\ 
\boxed{\Rea\rho=\tfrac12}.
\]

\section{本稿保留哪些标准基础}
下文使用 $\zeta$ 的解析延拓、在 $1$ 的 Laurent 展开、Euler 乘积和函数方程；还使用 Gamma 函数的零点与极点性质、Cauchy 积分公式、恒等定理，以及全纯函数的零点孤立性。这些是已建立的复分析与特殊函数基础，可分别在所附 Mathlib 源码和资料[4]--[5]中核对。

决定性的\emph{新增算术输入}是第\ref{chap:arithmetic}章的窗口定理。它不是从上述标准基础中直接推出的初等估计，本稿不会把它隐藏在“容易证明”之类的句子中。

\section{先用一个简单函数理解极点和高阶导数}
复数 $s=\beta+it$ 对应平面上的点 $(\beta,t)$；$\Rea s=\beta$ 是横坐标。临界线 $\Rea s=1/2$ 因而是一条竖直线。说零点在临界线上，就是说它们的横坐标全部等于 $1/2$，并没有说虚部也相同。

在一个区域内可作复微分的函数称为全纯函数。若某点附近可以写成
\[
f(s)=\frac{c}{s-\rho}+h(s),\qquad c\ne0,
\]
其中 $h$ 在该点附近全纯，则称 $\rho$ 为 $f$ 的单极点。$c/(s-\rho)$ 是趋近该点时不能被有限光滑部分消去的发散项。允许有孤立极点的函数称为亚纯函数。

先只看最简单的模型 $f(s)=c/(s-\rho)$。连续求导给出
\[
f^{(k)}(z)=\frac{(-1)^k k!c}{(z-\rho)^{k+1}},
\qquad
\left|\frac{(-1)^k}{k!}f^{(k)}(z)\right|
=\frac{|c|}{|z-\rho|^{k+1}}.
\]
因此，去掉阶乘以后，第 $k$ 阶导数的大小仍记住了观察点到极点的距离。若极点距离为 $d$，而其他部分在更大半径 $r>d$ 的圆盘内全纯，那么两部分在第 $k$ 阶的尺度分别为 $d^{-k}$ 和 $r^{-k}$。它们的比值中出现
\[
\left(\frac dr\right)^k\longrightarrow0.
\]
这就是第\ref{chap:window}章中极点项最终不能被其余部分抵消的原因。正式证明还要把多个极点分离，并同时控制实轴上的导数，不能只用这个模型替代全部论证。

\section{阅读时可以反复查阅的记号表}
\begin{longtable}{@{}p{.20\textwidth}p{.73\textwidth}@{}}
\toprule
记号 & 在证明中的作用\\
\midrule
\endhead
$\rho=\beta+it$ & 假想的 $\zeta$ 零点；$\beta$ 是需要限制的实部。\\
$G(x)$ & 由素数幂加权和得到的归一化实函数，负责承接算术估计。\\
$K,\varepsilon$ & $K$ 先固定；再取小的正数 $\varepsilon$，让指数接近 $1/2$。\\
$Y,[Y,Y^{3/2}]$ & 好窗口的起点和完整区间；窗口可以出现在任意远处。\\
$U=\log Y$ & 对数坐标下的好尺度；窗口变为 $[U,3U/2]$。\\
$\Phi(s)$ & $G$ 的 Mellin 变换，也是 $g(u)=G(e^u)$ 的 Laplace 变换。\\
$\theta$ & 暂时要排除的零点边界，最后选择在 $1/2$ 与假想零点之间。\\
$\sigma,k$ & 观察点的实部及导数阶数；本稿取 $336$ 与 $\lceil420U\rceil$。\\
$d,r,q$ & 极点距离、全纯估计半径及比值 $q=d/r<1$。\\
\bottomrule
\end{longtable}
